\documentclass[journal]{IEEEtran}  % Only keep this line
\usepackage[utf8]{inputenc}
\usepackage{hyperref}
\usepackage{cite}
\usepackage{graphicx}
\usepackage{amsmath, amssymb}

\title{Polymer Optical Fiber for Angle and Torque Measurements of a Series Elastic Actuator’s Spring}
\author{Eric Förster \and Patrick Förster}
\date{\today}

\begin{document}

\maketitle

\begin{abstract}
    Series elastic actuators (SEAs) can provide low output impedance, bandwidth close to the human movement, and direct measurement of torque through the spring deflection. These advantages enable the application of SEAs in human active orthosis and exoskeletons. However, conventional technologies to measure the spring deflection are bulky, inhibit natural pattern of movement, or are sensitive to misalignments. This paper presents the application of polymer optical fiber (POF) as a sensor to measure the spring deflection to overcome some of the issues of conventional technologies, since it is compact, lightweight, and has electromagnetic immunity. Furthermore, the spring is employed to validate a torque sensor based on POF stress-optic effects. Results show high linearity of both sensors and mean squared errors below the encoder resolution employed as reference on dynamic measurements.
\end{abstract}

\section{Introduction}
TexLab is a cross-platform implementation of the Language Server Protocol for the \LaTeX{} typesetting system. It aims to produce high-quality code completion results. The server may be used with any editor that implements the Language Server Protocol. It is written in Rust, a blazingly fast systems programming language.

\begin{IEEEkeywords}
    Language Server Protocol, Polymer optical fiber, Series elastic actuator, torque sensor, angle sensor.
\end{IEEEkeywords}

\section{Features}
The language server implements most of the Language Server Protocol specification. In addition to that, it implements additional functionality like building and forward search.

\section{Availability}
TexLab is available on \href{https://github.com/latex-lsp/texlab}{GitHub}, various package managers, and CTAN. Pre-compiled binaries are available on the \href{https://github.com/latex-lsp/texlab/releases}{GitHub Releases} page. Some editor extensions are able to automatically download TexLab.

\section{Installation}
There are various ways to install TexLab:
\begin{itemize}
    \item TexLab is included in some package managers like \texttt{brew}, \texttt{pacman}, and \texttt{scoop}.
    \item You can download a pre-compiled binary from our \href{https://github.com/latex-lsp/texlab/releases}{GitHub Releases} page.
    \item Some extensions like the Visual Studio Code extension or \texttt{coc-texlab} can automatically download the server for you.
    \item You can download the sources from either GitHub or CTAN and compile the server with \texttt{cargo build --release}.
\end{itemize}

\section{Usage}
\subsection{Synopsis}
\texttt{texlab [FLAGS] [OPTIONS]}

\subsection{Flags}
\begin{itemize}
    \item \texttt{-h}, \texttt{--help} Prints help information.
    \item \texttt{-q}, \texttt{--quiet} No output printed to stderr.
    \item \texttt{-V}, \texttt{--version} Prints version information.
    \item \texttt{-v}, \texttt{--verbosity} Increase message verbosity (\texttt{-vvvv} for max verbosity).
\end{itemize}

\section{Conclusion}
The study successfully demonstrates the use of polymer optical fibers for measuring angle and torque in series elastic actuators. The system is shown to be compact, lightweight, and reliable for robotic applications.

\end{document}
